%! Justifying a Claim Based on a Confidence Interval for a Difference of Population Proportions — Unit 6, Lesson 6.9 — Guided Notes (Answer Key)
\documentclass[10pt]{article}
\usepackage{apstats-article}
\usepackage{apstats-key}   % pulls in -boxes; do NOT also load -boxes
\newcommand{\TallMath}[1]{$\displaystyle #1\rule[-1.4em]{0pt}{3.2em}$}

\begin{document}

\pageheader{Unit 6, Lesson 6.9}{Guided Notes --- Answer Key}
\namedateperiod

\vspace{0.15in}

\begin{objectivebox}
\textbf{Lesson 6.9 Objectives:} Interpret a two-sample $z$-interval for $p_1-p_2$ and use it
to evaluate claims about the difference of two population proportions. \hfill \textit{UNC-4.M--N, Skills 3.D \& 4.B}
\end{objectivebox}

\begin{notesbox}{Interpreting a Two-Sample CI}
\textbf{Template:}\\
``We are \ans{95}\% confident that the true difference \ans{($p_F - p_M$)} in the proportions
of \ans{female students} and \ans{male students} who [context] is between \ans{$-0.049$} and \ans{$0.149$}.''\\[6pt]
\textbf{Example.} 95\% CI for $p_F - p_M$ is $(-0.049, 0.149)$:\\
\ansline{We are 95\% confident that the true difference ($p_F - p_M$) in the proportions of female and male students planning to take AP Chemistry is between $-0.049$ and $0.149$.}
\end{notesbox}

\begin{notesbox}{The Zero Rule}
\renewcommand{\arraystretch}{1.8}
\begin{tabular}{@{} p{0.44\linewidth} p{0.50\linewidth} @{}}
\textbf{0 is NOT in the interval} & \textbf{0 IS in the interval} \\
$\Rightarrow$ 0 is \ans{not plausible} & $\Rightarrow$ 0 is \ans{plausible} \\
$\Rightarrow$ \ans{convincing} evidence & $\Rightarrow$ \ans{no convincing} evidence \\
\quad of a difference & \quad of a difference \\
\end{tabular}\\[6pt]
\textbf{Sign:} If the entire interval is positive $\Rightarrow p_1$ \ans{$>$} $p_2$. \quad
Entirely negative $\Rightarrow p_1$ \ans{$<$} $p_2$.
\end{notesbox}

\begin{notesbox}{Worked Examples}
\textbf{Example 1.} A 90\% CI for $p_{\text{urban}} - p_{\text{rural}} = (0.04, 0.18)$.\\
Contains 0? \ans{No}. Conclusion: \ansline{Convincing evidence that urban residents have a higher proportion than rural residents.}

\textbf{Example 2.} A 95\% CI for $p_A - p_B = (-0.03, 0.09)$.\\
Contains 0? \ans{Yes}. Conclusion: \ansline{No convincing evidence of a difference between $p_A$ and $p_B$.}

\textbf{Example 3.} A 99\% CI for $p_{\text{new}} - p_{\text{old}} = (-0.22, -0.06)$.\\
Contains 0? \ans{No}. Conclusion: \ansline{Convincing evidence that the old method has a higher proportion than the new method ($p_{\text{old}} > p_{\text{new}}$).}
\end{notesbox}

\begin{practicebox}
The 95\% two-sample CI for the difference in the proportions of students who pass a
standardized test between schools A and B ($p_A - p_B$) is $(0.07, 0.23)$.

(a) \ansline{We are 95\% confident that the true difference ($p_A - p_B$) in the proportions of students at schools A and B who pass the standardized test is between 0.07 and 0.23.}

(b) \ansline{Yes. Since the entire interval $(0.07, 0.23)$ is positive (does not contain 0), there is convincing evidence that school A's pass rate is higher than school B's.}
\end{practicebox}

\end{document}
